% GENERATED FILE. Edit canonical lessons, then run npm run generate:books.
% canonical-source-hash: fnv1a64:95ad8daa950c721a
% canonical-lessons: PA-C353-latu, PA-C353-mukabla, PA-C353-kuddna, PA-C353-udauna, PA-C353-khidauna

\chapter{Out to Play}
\label{ch:pa-a2-out-to-play}
\hlchaptermodality{353}

\begin{chapteropening}
\textbf{By the end of this chapter:} \emph{I can say what children like to play outside.}

Complete the last lesson of chapter 353: I can say what children like to play outside.
\end{chapteropening}

\section[\texorpdfstring{\pa{ਲਾਟੂ} (lāṭū)}{ਲਾਟੂ (lāṭū)}]{\pa{ਲਾਟੂ} (lāṭū) — a spinning top}

\label{lesson:PA-C353-latu}



\begin{quote}
Before the new one: say the Punjabi for a game piece, then the Punjabi for a prize.
\end{quote}

\subsection*{You'll want to know: \pa{ਲਾਟੂ}}
\textbf{\pa{ਲਾਟੂ}} — \emph{lāṭū} — \textquotedblleft{}a spinning top\textquotedblright{}.

\begin{quote}
\pa{ਮੇਰੀ} \pa{ਛੋਟੀ} \pa{ਭੈਣ} \pa{ਕੋਲ} \pa{ਇੱਕ} \pa{ਲਾਲ} \pa{ਲਾਟੂ} \pa{ਹੈ।}
\end{quote}

\begin{quote}
\emph{merī chhoṭī bhaiṇ kol ikk lāl lāṭū hai.}
\end{quote}

\begin{quote}
My little sister has a red spinning top.
\end{quote}

A masculine noun.

\begin{cousinweb}
\textbf{\pa{ਲਾਟੂ}} (\emph{lāṭū}) is inherited from Sanskrit \emph{laṭati}, \textquotedblleft{}shakes\textquotedblright{} (a sense CDIAL infers); CDIAL lists Punjabi \emph{lāṭū}, \textquotedblleft{}spinning top\textquotedblright{}, under it, with a query (CDIAL 10916).
\end{cousinweb}

\begin{grammarlens}[title={What you've built}]
One more word for this chapter.
\end{grammarlens}

\subsection*{Your turn}
\begin{itemize}
\raggedright
  \item \emph{Say it:} \emph{lāṭū}
  \item \emph{Say it:} \emph{lāṭū}, once more
  \item \emph{Say it:} \emph{merī chhoṭī bhaiṇ kol ikk lāl lāṭū hai.} — My little sister has a red spinning top.
  \item \emph{Recall:} say the Punjabi for a game piece, then the Punjabi for a prize, then say \emph{lāṭū} again
\end{itemize}

\begin{tcolorbox}[breakable,colback=teal!4,colframe=teal!35!black,title={Before you move on}]
What does \pa{ਲਾਟੂ} mean? (\textquotedblleft{}A spinning top\textquotedblright{}.) Say it once more.
\end{tcolorbox}
\section[\texorpdfstring{\pa{ਮੁਕਾਬਲਾ} (mukāblā)}{ਮੁਕਾਬਲਾ (mukāblā)}]{\pa{ਮੁਕਾਬਲਾ} (mukāblā) — a contest; a competition}

\label{lesson:PA-C353-mukabla}



\begin{quote}
Before the new one: say the Punjabi for a prize, then the Punjabi for a spinning top.
\end{quote}

\subsection*{You'll want to know: \pa{ਮੁਕਾਬਲਾ}}
\textbf{\pa{ਮੁਕਾਬਲਾ}} — \emph{mukāblā} — \textquotedblleft{}a contest; a competition\textquotedblright{}.

\begin{quote}
\pa{ਅੱਜ} \pa{ਸਕੂਲ} \pa{ਵਿੱਚ} \pa{ਇੱਕ} \pa{ਮੁਕਾਬਲਾ} \pa{ਹੈ।}
\end{quote}

\begin{quote}
\emph{ajj sakūl vich ikk mukāblā hai.}
\end{quote}

\begin{quote}
There is a contest at school today.
\end{quote}

A masculine noun.

\begin{cousinweb}
\textbf{\pa{ਮੁਕਾਬਲਾ}} (\emph{mukāblā}) is a loan, usually traced through Persian to Arabic \emph{muqābala}, \textquotedblleft{}facing each other\textquotedblright{}.
\end{cousinweb}

\begin{grammarlens}[title={What you've built}]
One more word for this chapter.
\end{grammarlens}

\subsection*{Your turn}
\begin{itemize}
\raggedright
  \item \emph{Say it:} \emph{mukāblā}
  \item \emph{Say it:} \emph{mukāblā}, once more
  \item \emph{Say it:} \emph{ajj sakūl vich ikk mukāblā hai.} — There is a contest at school today.
  \item \emph{Recall:} say the Punjabi for a prize, then the Punjabi for a spinning top, then say \emph{mukāblā} again
  \item \emph{Write it:} copy \textbf{\pa{ਮੁਕਾਬਲਾ}} once, checking the whole word
\end{itemize}

\begin{tcolorbox}[breakable,colback=teal!4,colframe=teal!35!black,title={Before you move on}]
What does \pa{ਮੁਕਾਬਲਾ} mean? (\textquotedblleft{}A contest; a competition\textquotedblright{}.) Say it once more.
\end{tcolorbox}
\section[\texorpdfstring{\pa{ਕੁੱਦਣਾ} (kuddṇā)}{ਕੁੱਦਣਾ (kuddṇā)}]{\pa{ਕੁੱਦਣਾ} (kuddṇā) — to leap; to jump in}

\label{lesson:PA-C353-kuddna}



\begin{quote}
Before the new one: say the Punjabi for a spinning top, then the Punjabi for a contest.
\end{quote}

\subsection*{You'll want to know: \pa{ਕੁੱਦਣਾ}}
\textbf{\pa{ਕੁੱਦਣਾ}} — \emph{kuddṇā} — \textquotedblleft{}to leap; to jump in\textquotedblright{}.

\begin{quote}
\pa{ਬੱਚੇ} \pa{ਪਾਣੀ} \pa{ਵਿੱਚ} \pa{ਕੁੱਦਦੇ} \pa{ਹਨ} \pa{ਅਤੇ} \pa{ਤਰਦੇ} \pa{ਹਨ।}
\end{quote}

\begin{quote}
\emph{bachche pāṇī vich kuddde han ate tarde han.}
\end{quote}

\begin{quote}
The children jump into the water and swim.
\end{quote}

Past (I …): \pa{ਮੈਂ} \pa{ਪਾਣੀ} \pa{ਵਿੱਚ} \pa{ਕੁੱਦਿਆ}, \emph{maiṁ pāṇī vich kuddiā}.

\begin{cousinweb}
\textbf{\pa{ਕੁੱਦਣਾ}} (\emph{kuddṇā}) is inherited from Sanskrit \emph{kūrdati}, \textquotedblleft{}leaps, jumps\textquotedblright{}; CDIAL lists Punjabi \emph{kuddṇā}, \textquotedblleft{}to leap\textquotedblright{}, under it (CDIAL 3412).
\end{cousinweb}

\begin{grammarlens}[title={What you've built}]
One more word for this chapter.
\end{grammarlens}

\subsection*{Your turn}
\begin{itemize}
\raggedright
  \item \emph{Say it:} \emph{kuddṇā}
  \item \emph{Say it:} \emph{kuddṇā}, once more
  \item \emph{Say it:} \emph{bachche pāṇī vich kuddde han ate tarde han.} — The children jump into the water and swim.
  \item \emph{Recall:} say the Punjabi for a spinning top, then the Punjabi for a contest, then say \emph{kuddṇā} again
  \item \emph{Write it:} copy \textbf{\pa{ਕੁੱਦਣਾ}} once, checking the whole word
\end{itemize}

\begin{tcolorbox}[breakable,colback=teal!4,colframe=teal!35!black,title={Before you move on}]
What does \pa{ਕੁੱਦਣਾ} mean? (\textquotedblleft{}To leap; to jump in\textquotedblright{}.) Say it once more.
\end{tcolorbox}
\section[\texorpdfstring{\pa{ਉਡਾਉਣਾ} (uḍāuṇā)}{ਉਡਾਉਣਾ (uḍāuṇā)}]{\pa{ਉਡਾਉਣਾ} (uḍāuṇā) — to make something fly; to fly a kite}

\label{lesson:PA-C353-udauna}



\begin{quote}
Before the new one: say the Punjabi for a contest, then the Punjabi for to leap.
\end{quote}

\subsection*{You'll want to know: \pa{ਉਡਾਉਣਾ}}
\textbf{\pa{ਉਡਾਉਣਾ}} — \emph{uḍāuṇā} — \textquotedblleft{}to make something fly; to fly a kite\textquotedblright{}.

\begin{quote}
\pa{ਮੇਰਾ} \pa{ਦਾਦਾ} \pa{ਛੱਤ} \pa{ਉੱਤੇ} \pa{ਕਬੂਤਰ} \pa{ਉਡਾਉਂਦਾ} \pa{ਹੈ।}
\end{quote}

\begin{quote}
\emph{merā dādā chhatt utte kabūtar uḍāundā hai.}
\end{quote}

\begin{quote}
My grandfather flies pigeons from the roof.
\end{quote}

Past (I …): \pa{ਮੈਂ} \pa{ਕਬੂਤਰ} \pa{ਉਡਾਏ}, \emph{maiṁ kabūtar uḍāe}.

\begin{cousinweb}
\textbf{\pa{ਉਡਾਉਣਾ}} (\emph{uḍāuṇā}) is inherited from Sanskrit \emph{uḍḍīyate}, \textquotedblleft{}flies up\textquotedblright{}; CDIAL lists Punjabi \emph{uḍāuṇā}, \textquotedblleft{}to cause to fly\textquotedblright{}, under it (CDIAL 1697).
\end{cousinweb}

\begin{grammarlens}[title={What you've built}]
One more word for this chapter.
\end{grammarlens}

\subsection*{Your turn}
\begin{itemize}
\raggedright
  \item \emph{Say it:} \emph{uḍāuṇā}
  \item \emph{Say it:} \emph{uḍāuṇā}, once more
  \item \emph{Say it:} \emph{merā dādā chhatt utte kabūtar uḍāundā hai.} — My grandfather flies pigeons from the roof.
  \item \emph{Recall:} say the Punjabi for a contest, then the Punjabi for to leap, then say \emph{uḍāuṇā} again
  \item \emph{Write it:} copy \textbf{\pa{ਉਡਾਉਣਾ}} once, checking the whole word
\end{itemize}

\begin{tcolorbox}[breakable,colback=teal!4,colframe=teal!35!black,title={Before you move on}]
What does \pa{ਉਡਾਉਣਾ} mean? (\textquotedblleft{}To make something fly; to fly a kite\textquotedblright{}.) Say it once more.
\end{tcolorbox}
\section[\texorpdfstring{\pa{ਖਿਡਾਉਣਾ} (khiḍāuṇā)}{ਖਿਡਾਉਣਾ (khiḍāuṇā)}]{\pa{ਖਿਡਾਉਣਾ} (khiḍāuṇā) — to play with a child; to keep a child amused}

\label{lesson:PA-C353-khidauna}



\begin{quote}
Before the new one: say the Punjabi for to leap, then the Punjabi for to make something fly.
\end{quote}

\subsection*{You'll want to know: \pa{ਖਿਡਾਉਣਾ}}
\textbf{\pa{ਖਿਡਾਉਣਾ}} — \emph{khiḍāuṇā} — \textquotedblleft{}to play with a child; to keep a child amused\textquotedblright{}.

\begin{quote}
\pa{ਮੇਰੀ} \pa{ਭੈਣ} \pa{ਬੱਚੇ} \pa{ਨੂੰ} \pa{ਵਿਹੜੇ} \pa{ਵਿੱਚ} \pa{ਖਿਡਾਉਂਦੀ} \pa{ਹੈ।}
\end{quote}

\begin{quote}
\emph{merī bhaiṇ bachche nū̃ vihṛe vich khiḍāundī hai.}
\end{quote}

\begin{quote}
My sister plays with the baby in the courtyard.
\end{quote}

Past (I …): \pa{ਮੈਂ} \pa{ਬੱਚੇ} \pa{ਨੂੰ} \pa{ਖਿਡਾਇਆ}, \emph{maiṁ bachche nū̃ khiḍāiā}.

\begin{cousinweb}
\textbf{\pa{ਖਿਡਾਉਣਾ}} (\emph{khiḍāuṇā}) is inherited from a reconstructed \emph{*khēḍ}, \textquotedblleft{}play\textquotedblright{}; CDIAL lists Punjabi \emph{khiḍāuṇā}, \textquotedblleft{}play\textquotedblright{}, under it (CDIAL 3918).
\end{cousinweb}

\begin{grammarlens}[title={What you've built}]
One more word for this chapter.
\end{grammarlens}

\subsection*{Your turn}
\begin{itemize}
\raggedright
  \item \emph{Say it:} \emph{khiḍāuṇā}
  \item \emph{Say it:} \emph{khiḍāuṇā}, once more
  \item \emph{Say it:} \emph{merī bhaiṇ bachche nū̃ vihṛe vich khiḍāundī hai.} — My sister plays with the baby in the courtyard.
  \item \emph{Recall:} say the Punjabi for to leap, then the Punjabi for to make something fly, then say \emph{khiḍāuṇā} again
  \item \emph{Write it:} copy \textbf{\pa{ਖਿਡਾਉਣਾ}} once, checking the whole word
\end{itemize}

\begin{tcolorbox}[breakable,colback=teal!4,colframe=teal!35!black,title={Before you move on}]
What does \pa{ਖਿਡਾਉਣਾ} mean? (\textquotedblleft{}To play with a child; to keep a child amused\textquotedblright{}.) Say it once more.
\end{tcolorbox}
